% Laporan temuan EDA kelompok, memakai huruf dan tata halaman template FILKOM UB
% (../template-latex), tetapi tanpa halaman daftar dan tanpa halaman bab: bentuknya
% laporan tugas biasa dengan bagian bernomor 1, 2, 3, ...
% Kompilasi: latexmk main.tex   (pdfLaTeX + biber)
\documentclass[12pt,a4paper,oneside]{report}
\usepackage{gaya/ubtesis}
\usepackage{makecell}  % sel tabel bertumpuk
\addbibresource{referensi.bib}

% Penomoran laporan: 1, 1.1, Tabel 1, Gambar 1, (1) — tidak mengikuti nomor bab.
\renewcommand\thesection{\arabic{section}}
\counterwithout{table}{chapter}
\counterwithout{figure}{chapter}
\counterwithout{equation}{chapter}
\renewcommand\theequation{\arabic{equation}}
\defbibheading{laporan}{\section*{Daftar Referensi}}
% Judul bagian utama tidak boleh tertinggal di kaki halaman tanpa isinya.
\pretocmd{\section}{\needspace{12\baselineskip}}{}{}

% ------------------------------------------------------------- data sampul
\newcommand\judul{Laporan Temuan EDA: Kurasi Dataset dan Karakteristik
  Modalitas Tobacco3482}
\newcommand\matakuliah{ANALISIS \& PEMROSESAN DATA MULTIMODAL}
\newcommand\pengampu{Dr. Eng. Fitra Abdurrachman Bachtiar, S.T., M.Eng.}
\newcommand\programstudi{PROGRAM STUDI MAGISTER ILMU KOMPUTER}
\newcommand\departemen{DEPARTEMEN TEKNIK INFORMATIKA}
\newcommand\fakultas{FAKULTAS ILMU KOMPUTER}
\newcommand\universitas{UNIVERSITAS BRAWIJAYA}
% Anggota kelompok, diurutkan menurut NIM.
\newcommand\anggota{%
  Yulian Anggreaif Rudamaga & 256150100011001 \\
  Zalzabila Rani             & 256150100011008 \\
  Ramadhani Al-Qadri         & 256150100011010 \\}

\begin{document}
\margininduk
% Sampul, mengikuti susunan sampul Progress 1 kelompok.
\begin{titlepage}
\thispagestyle{empty}
\setlength{\parskip}{0pt}
\centering
\vspace*{0.4cm}
{\bfseries\ukuran{14}\matakuliah\par}
\vspace{0.8cm}
{\bfseries\ukuran{16}\judul\par}
\vspace{1.8cm}
\includegraphics[width=4.5cm]{logo-ub}\par
\vspace{1.8cm}
{\ukuran{12}Disusun oleh:\par\vspace{0.25cm}
 \begin{tabular}{@{}l@{\quad:\quad}l@{}}\anggota\end{tabular}\par}
\vspace{1.2cm}
{\ukuran{12}Dosen Pengampu:\par
 \pengampu\par}
\vfill
{\bfseries\ukuran{14}\programstudi\par
 \departemen\par
 \fakultas\par
 \universitas\par}
\vspace{0.5cm}
\end{titlepage}
\setcounter{page}{1}


\marginisi
\pagenumbering{arabic}
\section{Pendahuluan}\label{sec:pendahuluan}

Laporan ini merangkum hasil kurasi dataset dan \textit{exploratory data analysis} (EDA)
atas Tobacco3482 \parencite{kumar2014}. Dataset ini terdiri atas dua modalitas:

\begin{itemize}
  \item \textbf{citra halaman}, berupa versi JPG dari Kaggle \parencite{kaggletobacco};
  \item \textbf{teks OCR} dari QS-OCR \parencite{qsocr}, hasil Tesseract
  \parencite{smith2007}.
\end{itemize}

Seluruh angka berasal dari eksekusi notebook \texttt{11\_eda.ipynb}. Notebook itu tidak
melatih model; fitur diambil dari model yang dibekukan.

Analisis yang memakai label hanya memakai bagian latih seed~42 (800 dokumen). Analisis
ini mencakup kata khas per kelas, kemiripan antar-kelas, dan classifier kode ID. Data uji
hanya dipakai untuk mengukur akurasi \textit{probe}. Dengan cara ini EDA tidak
\enquote{mengintip} data uji yang nanti dipakai eksperimen.

Perbandingan antar-mesin OCR belum termasuk di sini. Bagian itu menyusul setelah OCR
multi-mesin dijalankan.

\section{Dataset dan Kurasi}\label{sec:kurasi}

\subsection{Pembagian Data}

Ke-3.482 pasangan citra--teks lengkap; tidak ada citra atau teks yang hilang. Kelasnya
timpang, dengan rasio kelas terbesar (Memo, 620) berbanding terkecil (Resume, 120)
sebesar 5,17.

Pembagian data mengikuti \textcite{audebert2020}: 800 dokumen latih dan sisanya uji.
Setiap pembagian terstratifikasi per kelas, untuk tiga seed (42, 43, 44). Komposisinya
identik di ketiga seed (Tabel~\ref{tab:split}); yang berbeda hanya dokumen yang terpilih.

\begin{table}[H]
  \centering
  \caption{Jumlah dokumen per kelas pada setiap bagian (sama untuk seed 42, 43, 44)}
  \label{tab:split}
  \fontsize{10}{12}\selectfont
  \begin{tabular}{|l|c|c|c|c|c|c|c|c|c|c|c|}
    \hline
    & \rotatebox{90}{Advertisement} & \rotatebox{90}{Email} & \rotatebox{90}{Form} &
    \rotatebox{90}{Letter} & \rotatebox{90}{Memo} & \rotatebox{90}{News} &
    \rotatebox{90}{Note} & \rotatebox{90}{Report} & \rotatebox{90}{Resume} &
    \rotatebox{90}{Scientific} & \textbf{Total} \\ \hline
    Train & 48 & 124 & 89 & 117 & 128 & 39 & 41 & 55 & 25 & 54 & 720 \\ \hline
    Val   & 5 & 14 & 10 & 13 & 14 & 4 & 5 & 6 & 3 & 6 & 80 \\ \hline
    Test  & 177 & 461 & 332 & 437 & 478 & 145 & 155 & 204 & 92 & 201 & 2.682 \\ \hline
  \end{tabular}
\end{table}

\subsection{Aturan Kurasi}

Kurasi dilakukan dengan \textbf{menandai} dokumen, bukan membuangnya. Derau OCR adalah
objek penelitian, sehingga dokumen yang bising tetap dipakai dan hanya diberi penanda.

Ambang setiap penanda diambil dari ekor distribusi: 5\% terbawah, atau 5\% teratas untuk
OOV. Ambang ini tidak ditetapkan sebelum melihat data. Hasilnya disajikan pada
Tabel~\ref{tab:penanda}.

\begin{table}[H]
  \centering
  \caption{Penanda kurasi, ambang yang diperoleh dari data, dan jumlah dokumen tertandai}
  \label{tab:penanda}
  \fontsize{10}{12}\selectfont
  \begin{tabularx}{\textwidth}{|l|>{\raggedright\arraybackslash}X|r|}
    \hline
    \textbf{Penanda} & \textbf{Aturan (ambang dari data)} & \textbf{Dokumen} \\ \hline
    Teks kosong & tidak ada kata alfabetik sama sekali & 37 \\ \hline
    Teks pendek & $\leq$ 23 kata alfabetik (di antara yang tidak kosong) & 163 \\ \hline
    OOV tinggi & OOV $\geq$ 0,369 & 172 \\ \hline
    Resolusi rendah & perkiraan $<$ 203 dpi & 68 \\ \hline
    Ketajaman rendah & variansi Laplacian $<$ 788 & 175 \\ \hline
    Duplikat & korelasi \textit{thumbnail} $\geq$ 0,98 dan teks mirip & 9 \\ \hline
  \end{tabularx}
\end{table}

Duplikat diperiksa dalam dua tahap. Tahap pertama mencari pasangan halaman yang
\textit{thumbnail}-nya sangat mirip. Tahap kedua mengonfirmasi pasangan itu lewat
kemiripan teks OCR-nya. Dua tahap diperlukan karena halaman dokumen sebagian besar putih:
kemiripan citra saja akan menyamakan formulir yang isinya berbeda tetapi templatnya sama.

Hasilnya 12 pasangan yang melibatkan 9 dokumen. Pasangan yang melintasi bagian latih dan
uji, yaitu kebocoran data, ada 1 pada seed~42, 5 pada seed~43, dan 4 pada seed~44.

Ada satu catatan. Enam dari sembilan dokumen duplikat berasal dari kelas Advertisement,
kelas yang juga paling banyak teks kosongnya. Pasangan yang \emph{keduanya} tanpa teks
dikonfirmasi otomatis, karena tidak ada teks untuk dibandingkan. Karena itu duplikat ini
perlu diperiksa secara visual sebelum dipakai untuk mengubah pembagian data.

Hasil kurasi disimpan sebagai \texttt{manifest\_kurasi.csv}. Isinya satu baris per
dokumen: label, kolom pembagian untuk ketiga seed, seluruh statistik dan penanda, serta
penanda kebocoran. Berkas citra dan teks sendiri tidak disalin, karena lisensi
Tobacco3482 tidak memperbolehkan redistribusi.

\section{Modalitas Citra}\label{sec:citra}

\subsection{Statistik Halaman}

Citra JPG dari Kaggle beresolusi tinggi:
\begin{itemize}
  \item lebar median 1.728 piksel, sekitar 203 dpi untuk kertas \textit{letter};
  \item kuartil atas 2.544 piksel, sekitar 299 dpi;
  \item rentangnya 96 sampai 414 dpi.
\end{itemize}

Statistiknya per kelas ditunjukkan Gambar~\ref{fig:citra}.

\begin{figure}[H]
  \centering
  \includegraphics[width=\textwidth]{eda_statistik_citra}
  \caption{Statistik citra per kelas: perkiraan DPI, kontras, rasio tinta, dan
  ketajaman (skala log)}
  \label{fig:citra}
\end{figure}

Gambar~\ref{fig:citra} memperlihatkan tiga hal.
\begin{enumerate}
  \item \textbf{Resolusi tidak seragam antar-kelas.} Pada Email dan Resume, median dan
  kedua kuartilnya sama-sama 203 dpi. Kelas lain umumnya menyebar antara 203 dan 300 dpi.
  Resolusi asli memang hilang setelah citra diubah ke $384\times384$, tetapi pola ini
  menunjukkan bahwa kelas-kelas tersebut dipindai dengan cara berbeda. Ini sumber potensi
  bias.
  \item \textbf{Advertisement jauh lebih gelap.} Kontras dan rasio tintanya paling
  tinggi, karena halaman iklan didominasi gambar, bukan tulisan.
  \item \textbf{Email dan Note paling jarang tintanya}, sehingga ketajamannya juga paling
  rendah. Ukuran ketajaman (variansi Laplacian) bergantung pada jumlah tepi. Halaman yang
  hanya berisi sedikit tulisan otomatis mendapat skor rendah walaupun tidak kabur.
  Karena itu penanda \enquote{ketajaman rendah} harus dibaca sebagai \enquote{sedikit
  tulisan atau kabur}, bukan semata \enquote{kabur}.
\end{enumerate}

\subsection{Fitur CNN}

Fitur 1.280 dimensi diambil dari MobileNetV2 berbobot ImageNet tanpa
\textit{fine-tune} \parencite{sandler2018}. Regresi logistik di atas fitur beku ini
(\textit{linear probe}) mencapai \textbf{OA 0,7293} dan macro F1 0,7099. Sebagai
pembanding, menebak kelas mayoritas hanya mencapai 0,1782. Jadi tata letak halaman sudah
sangat informatif bahkan tanpa pelatihan khusus dokumen.

Pada proyeksi t-SNE (Gambar~\ref{fig:tsne}, kiri), Email, Advertisement, Form, News,
Resume, dan Scientific membentuk kelompok yang jelas. Letter, Memo, dan Report bercampur
di tengah. Ketiganya adalah dokumen ketikan berparagraf yang tata letaknya mirip.

\section{Modalitas Teks}\label{sec:teks}

\subsection{Statistik Teks}

\begin{itemize}
  \item \textbf{Panjang:} median 174 kata per dokumen, rata-rata 224, maksimum 1.579.
  \item \textbf{Teks kosong:} 37 dokumen tidak menghasilkan satu pun kata alfabetik.
  \item \textbf{OOV:} median 5,5\%, rata-rata 9,8\%. OOV di sini adalah bagian kata
  alfabetik minimal empat huruf yang tidak ada dalam daftar kata bahasa Inggris
  \texttt{wordfreq} \parencite{wordfreq}.
\end{itemize}

Daftar kata itu longgar: salah baca yang kebetulan berupa kata yang ada, misalnya
\texttt{tbe} untuk \texttt{the}, tidak terhitung. Jadi angka OOV ini adalah \emph{batas
bawah} galat OCR. Audebert et al.\ melaporkan sekitar 26\%, tetapi dengan kamus yang
lebih ketat, sehingga kedua angka tidak dapat dibandingkan langsung.

\begin{figure}[H]
  \centering
  \includegraphics[width=\textwidth]{eda_statistik_teks}
  \caption{Jumlah kata dan OOV per dokumen menurut kelas}
  \label{fig:teks}
\end{figure}

\subsection{Kata Khas per Kelas}

Kata dengan bobot TF-IDF tertinggi per kelas (Tabel~\ref{tab:tfidf}) menunjukkan bahwa
teks OCR membawa penanda kelas yang kuat. Contohnya:
\begin{itemize}
  \item peringatan \textit{surgeon general} pada Advertisement;
  \item kepala surel pada Email;
  \item istilah akademik pada Resume dan Scientific.
\end{itemize}

Kelas Note berbeda. Daftarnya berisi potongan kata seperti \texttt{ore}, \texttt{ous},
dan \texttt{vay}, yang menandakan teks hasil OCR dari tulisan tangan atau catatan pendek
yang rusak.

\begin{table}[H]
  \centering
  \caption{Kata TF-IDF tertinggi per kelas (dihitung dari bagian latih seed 42)}
  \label{tab:tfidf}
  \fontsize{10}{12}\selectfont
  \begin{tabularx}{\textwidth}{|l|>{\raggedright\arraybackslash}X|}
    \hline
    \textbf{Kelas} & \textbf{Kata khas} \\ \hline
    Advertisement & warning, surgeon, cigarette, general, advertisement, smoking, tar, nicotine \\ \hline
    Email & sent, subject, message, original, wednesday, date, thanks, meeting \\ \hline
    Form & date, fax, number, pages, form, total, transmission, telephone \\ \hline
    Letter & tobacco, dear, research, new, sincerely, york, product, council \\ \hline
    Memo & correspondence, attached, memorandum, inter, office, subject, date \\ \hline
    News & said, smoking, health, smoke, tobacco, smokers, cancer, industry \\ \hline
    Note & desk, information, james, ore, web, ous, christine, ran, vay, site \\ \hline
    Report & tobacco, smoking, cigarette, products, smokers, study, test, research \\ \hline
    Resume & university, research, biology, assistant, professor, medicine, publications \\ \hline
    Scientific & cell, cells, research, results, study, induced, university, data \\ \hline
  \end{tabularx}
\end{table}

\subsection{Fitur Embedding Teks}

Setiap dokumen direpresentasikan sebagai rata-rata vektor FastText kata-katanya
\parencite{bojanowski2017}. \textit{Linear probe} di atas representasi ini mencapai
\textbf{OA 0,6685} dan macro F1 0,6491. Angka ini lebih rendah daripada citra (0,7293),
tetapi jauh di atas kelas mayoritas.

Pada t-SNE (Gambar~\ref{fig:tsne}, kanan), Resume dan Email membentuk kelompok yang
sangat jelas, Scientific mengelompok di bagian atas, dan Form di sisi kanan. Pola
pengelompokannya berbeda dari citra. Perbedaan inilah yang menjadi dasar harapan bahwa
kedua modalitas saling melengkapi.

\begin{figure}[H]
  \centering
  \includegraphics[width=\textwidth]{eda_tsne_citra_teks}
  \caption{Proyeksi t-SNE fitur citra (kiri) dan fitur teks (kanan), diwarnai menurut
  kelas}
  \label{fig:tsne}
\end{figure}

\section{Komplementaritas Kedua Modalitas}\label{sec:komplementaritas}

Gambar~\ref{fig:komplementaritas} membandingkan F1 \textit{probe} per kelas antara teks
dan citra.

\begin{figure}[H]
  \centering
  \includegraphics[width=0.8\textwidth]{eda_komplementaritas}
  \caption{Selisih F1 \textit{probe} teks dikurangi citra per kelas}
  \label{fig:komplementaritas}
\end{figure}

\begin{itemize}
  \item \textbf{Teks lebih unggul} pada Resume (+0,237), Scientific (+0,076), dan Email
  (+0,018). Resume adalah contoh terkuat: tata letaknya beragam, tetapi kosakatanya
  sangat khas.
  \item \textbf{Citra lebih unggul} pada Advertisement ($-$0,297) dan Note ($-$0,289).
  Keduanya adalah kelas yang teksnya paling sering hilang atau rusak
  (Subbagian~\ref{sec:bising}).
  \item \textbf{Selisihnya kecil} ($<$ 0,11) untuk Letter, Form, News, Report, dan Memo.
\end{itemize}

Ini \textit{probe} pada fitur beku dengan 800 dokumen latih, bukan model akhir. Namun
arahnya memberi gambaran awal di kelas mana modalitas teks diharapkan menyumbang paling
banyak.

\section{Modalitas Hilang dan Bising}\label{sec:bising}

Persentase dokumen tertandai per kelas disajikan pada Tabel~\ref{tab:penanda-kelas}.

\begin{table}[H]
  \centering
  \caption{Persentase dokumen tertandai per kelas}
  \label{tab:penanda-kelas}
  \fontsize{10}{12}\selectfont
  \begin{tabular}{|l|c|c|c|c|c|c|}
    \hline
    \textbf{Kelas} & \makecell{\textbf{Teks}\\\textbf{kosong}} &
    \makecell{\textbf{Teks}\\\textbf{pendek}} & \makecell{\textbf{OOV}\\\textbf{tinggi}} &
    \makecell{\textbf{Resolusi}\\\textbf{rendah}} &
    \makecell{\textbf{Ketajaman}\\\textbf{rendah}} & \textbf{Duplikat} \\ \hline
    Advertisement & 3,5 & 17,4 & 13,0 & 3,9 & 2,2 & 2,6 \\ \hline
    Email & 0,0 & 4,2 & 0,2 & 0,0 & 19,9 & 0,3 \\ \hline
    Form & 0,0 & 1,4 & 10,7 & 0,0 & 0,0 & 0,0 \\ \hline
    Letter & 0,0 & 0,2 & 0,2 & 0,0 & 0,0 & 0,0 \\ \hline
    Memo & 0,0 & 0,2 & 0,6 & 6,5 & 0,0 & 0,0 \\ \hline
    News & 0,0 & 2,1 & 8,0 & 3,7 & 0,5 & 0,5 \\ \hline
    Note & 14,4 & 41,3 & 30,8 & 2,0 & 24,9 & 0,0 \\ \hline
    Report & 0,0 & 0,0 & 1,5 & 1,9 & 0,0 & 0,0 \\ \hline
    Resume & 0,0 & 0,0 & 0,0 & 1,7 & 0,0 & 0,0 \\ \hline
    Scientific & 0,0 & 1,1 & 3,4 & 0,4 & 0,0 & 0,0 \\ \hline
    \textbf{Semua} & 1,1 & 4,7 & 4,9 & 2,0 & 5,0 & 0,3 \\ \hline
  \end{tabular}
\end{table}

\textbf{Modalitas teks paling bermasalah pada Note.} Sebanyak 14,4\% dokumennya tanpa
teks sama sekali, 41,3\% teksnya pendek, dan 30,8\% OOV-nya tinggi. Note sering berupa
catatan tulisan tangan, yang memang di luar kemampuan OCR cetak.

Advertisement menyusul, dengan 17,4\% teks pendek dan 13,0\% OOV tinggi. Form juga
banyak ber-OOV tinggi (10,7\%), karena isinya angka, kode, dan singkatan.

\textbf{Mengapa iklan kehilangan teksnya.} Gambar~\ref{fig:contoh} memperlihatkan
dokumen yang teksnya paling bermasalah. Beberapa iklan dipindai dalam posisi
\textbf{terputar 90°}. Tesseract pada konfigurasi QS-OCR (\texttt{-{}-psm 3}) tidak
mendeteksi orientasi, sehingga teks pada halaman terputar tidak terbaca. Temuan ini
relevan untuk perbandingan mesin OCR berikutnya: mesin dengan koreksi orientasi dapat
berbeda jauh justru pada kelas ini.

\begin{figure}[H]
  \centering
  \includegraphics[width=\textwidth]{eda_contoh_teks_bermasalah}
  \caption{Contoh dokumen dengan modalitas teks paling bermasalah}
  \label{fig:contoh}
\end{figure}

\textbf{Dokumen yang hanya satu modalitasnya bermasalah} adalah tempat fusion seharusnya
paling menolong. Teks bermasalah dengan citra baik paling banyak ditemukan pada Note
(101 dokumen), Advertisement (62), dan Form (48). Citra bermasalah dengan teks baik
paling banyak pada Email (94) dan Memo (38).

Angka Email perlu dibaca hati-hati. Seluruhnya berasal dari penanda ketajaman
rendah, yang pada Email lebih mencerminkan sedikitnya tulisan daripada citra yang kabur
(Subbagian~\ref{sec:citra}).

\section{Jalan Pintas Kode ID Arsip}\label{sec:kode-id}

\textcite{larson2025} menemukan bahwa kode ID arsip yang tercetak di halaman Tobacco3482
sudah cukup untuk memprediksi kelas sekitar 60\%. Kode semacam itu ikut terbaca OCR.
Dengan deteksi heuristik (deret minimal enam digit), kode ID ditemukan dalam teks OCR
pada 20\% (Memo) sampai 56\% (Email) dokumen.

Ada dua uji jalan pintas, keduanya memakai classifier yang hanya diberi tiga digit awal
kode:
\begin{itemize}
  \item \textbf{dari kode ID dalam teks OCR:} OA 0,2782, hanya 10 poin di atas kelas
  mayoritas (0,1782);
  \item \textbf{dari tiga karakter awal ID berkas:} OA \textbf{0,4720}.
\end{itemize}

Jadi sistem penomoran arsip memang membawa informasi kelas. Model penelitian ini tidak
melihat nama berkas, sehingga jalan pintas kedua tidak berlaku untuk model kita. Jalan
pintas pertama tampak lemah, tetapi heuristik di sini lebih kasar daripada anotasi
manual Larson et al., sehingga belum dapat dinyatakan tidak ada. Kontrol kode ID tetap
diperlukan, terutama pada tahap yang memakai posisi teks.

\section{Implikasi untuk Eksperimen}\label{sec:implikasi}

\begin{enumerate}
  \item \textbf{Pembagian data dipertahankan.} Hanya 1 sampai 5 pasangan duplikat
  melintasi latih dan uji per seed. Pasangan ini akan diperiksa visual dulu, lalu
  dilaporkan atau dikeluarkan dari data uji, bukan dengan membuat pembagian baru.
  \item \textbf{Kelas yang diharapkan paling diuntungkan teks adalah Resume, Scientific,
  dan Email.} Bila model gabungan tidak membaik di kelas-kelas ini, itu tanda bahwa cabang
  teks tidak dipakai.
  \item \textbf{Note dan Advertisement adalah titik lemah teks.} Keduanya juga kelas
  tempat mesin OCR lain, terutama yang mengoreksi orientasi, paling mungkin membuat
  perbedaan. Hasil per kelas harus dilaporkan, tidak hanya akurasi total.
  \item \textbf{Penanda kurasi dapat dipakai untuk analisis per kelompok,} misalnya
  akurasi model pada dokumen bertanda teks bermasalah dibandingkan dokumen lain.
  \item \textbf{Bias resolusi antar-kelas} dicatat sebagai keterbatasan dataset.
\end{enumerate}


\printbibliography[heading=laporan]
\end{document}
